%% ============================================================================
%%
%%   S U P E R S E D E D   --   D O   N O T   C I T E
%%
%% This manuscript has been withdrawn. It is retained in the repository so that
%% the correction which replaced it remains traceable, not because any claim in
%% it still stands.
%%
%% WHAT IN THIS MANUSCRIPT DID NOT SURVIVE MEASUREMENT
%%
%%   The attribution of per-request cost to KEY CONVERSION. This document
%%   reports that the dominant term is "key material re-converted from a string
%%   on every call" (see the abstract, the introduction, and contribution 2).
%%   Direct measurement of that conversion puts it at 0.583 us -- far too small
%%   to be the explanation. The dominant term is a discarded exception: the
%%   library attempts an ASYMMETRIC key parse before the symmetric one and
%%   throws the failure away on every call, at 18.04 us, 60% of the call.
%%
%%   Its explanation of the HMAC-versus-RSA asymmetry is backwards. It ascribes
%%   the difference to secret conversion being dearer than parsing a PEM; a
%%   successful asymmetric parse costs 13.83 us against 0.583 us for the
%%   conversion. The real asymmetry is that the RSA path's parse SUCCEEDS,
%%   making its cost useful work, while the HMAC path's FAILS.
%%
%%   Its in-situ versus isolated comparison attributes a difference to
%%   deployment across two environments that also differed by two OpenSSL major
%%   series, so the comparison cannot support that attribution.
%%
%%   Its stated runtime version for the system under test was read from the
%%   host, not from the container, and is wrong.
%%
%% A SEPARATE WITHDRAWN CLAIM, NOT MADE IN THIS FILE
%%
%%   An earlier manuscript from this project argued that layered Zero Trust
%%   enforcement drives the service toward EVENT-LOOP SATURATION, and proposed
%%   relocating verification into the Envoy sidecar before that point. Direct
%%   in-situ measurement did not reproduce that behaviour and the claim is
%%   withdrawn. That manuscript's source was never committed to this
%%   repository; this file does not contain the claim and in fact disclaims it
%%   (see its threats to validity: "We make no claim about throughput ceilings,
%%   saturation points or autoscaling thresholds"). It is named here only so a
%%   reader who arrives looking for it knows its status.
%%
%% THE CURRENT MANUSCRIPT
%%
%%   key-resolution-cost.tex, in this directory.
%%
%% WHERE THE CORRECTION LANDED
%%
%%   commit 0b8f08f620977c9c0fe8b0bfe60a0d2d70306602 (2026-09-18)
%%   "Phase 5: rewrite the manuscript around the measured mechanism"
%%   released in tag v2.0.0
%%
%% See also SUPERSEDED.md in this directory.
%%
%% ============================================================================

%% not-the-cryptography.tex
%%
%% "It Is Not the Cryptography: Where the Cost of Token Validation Actually
%% Goes in Cloud-Native Microservices"
%%
%% EVERY numeric value in this document is a macro from paper_macros.tex, which
%% is generated by analysis/make_paper_macros.py from files under data/runs/.
%% No number is typed inline. paper_macros_provenance.csv maps each macro to the
%% file it was read from.
\documentclass[pdflatex,sn-basic,twocolumn]{sn-jnl}

% sn-jnl.cls calls \SetFootnoteHook and \DeclareNewFootnote at
% \AtBeginDocument but does not itself load the package that defines them.
% Loading it here keeps the class's own footnote setup working.
% Scalable Latin Modern replaces the bitmap-era Computer Modern the class asks
% for, which XeTeX cannot supply at the class's fractional sizes.
\usepackage{lmodern}
\usepackage{manyfoot}
% The class also assumes amsmath is present (it calls \allowdisplaybreaks).
\usepackage{amsmath}
\usepackage{amssymb}
\usepackage{graphicx}
\usepackage{booktabs}
\usepackage{url}
\usepackage{xcolor}

% paper_macros.tex -- GENERATED. Do not edit by hand.
% Regenerate: .venv/bin/python analysis/make_paper_macros.py
% Every macro below is read out of a file under data/runs/; the
% macro -> file map is in paper_macros_provenance.csv.

\providecommand{\PLACEHOLDER}[1]{\textcolor{red}{[MISSING: #1]}}

\newcommand{\AbCpuDelta}{76.92}% data/runs/INSITU-AB-authmode-jwt/openloop_ramp.csv
\newcommand{\AbCpuJwt}{146.38}% data/runs/INSITU-AB-authmode-jwt/openloop_ramp.csv
\newcommand{\AbCpuNone}{69.45}% data/runs/INSITU-AB-authmode-none/openloop_ramp.csv
\newcommand{\AbCpuPerRequestUs}{384.6}% data/runs/INSITU-AB-authmode-jwt/openloop_ramp.csv
\newcommand{\AbLatencyJwt}{2.0332}% data/runs/INSITU-AB-authmode-jwt/openloop_ramp.csv
\newcommand{\AbLatencyNone}{1.6802}% data/runs/INSITU-AB-authmode-none/openloop_ramp.csv
\newcommand{\AbRate}{200}% data/runs/INSITU-AB-authmode-jwt/openloop_ramp.csv
\newcommand{\EnvKubernetesVersion}{1.32.2}% data/runs/2026-09-14T15-01-36Z-verifyrate-hs256-preparsed/run_metadata.json
\newcommand{\EnvNodeVersion}{26.6.0}% data/runs/2026-09-14T15-01-36Z-verifyrate-hs256-preparsed/run_metadata.json
\newcommand{\Fifty}{50}% data/runs/2026-09-14T15-01-36Z-verifyrate-hs256-preparsed/run_metadata.json
\newcommand{\Five}{5}% data/runs/2026-09-14T15-01-36Z-verifyrate-hs256-preparsed/run_metadata.json
\newcommand{\FourHundred}{400}% data/runs/2026-09-14T15-01-36Z-verifyrate-hs256-preparsed/run_metadata.json
\newcommand{\GapEffectPreparsedRatio}{24.2}% data/runs/MECHANISM-2026-09-14-interarrival-gap/gap_keyform_incontainer.csv
\newcommand{\GapEffectStringRatio}{2.7}% data/runs/MECHANISM-2026-09-14-interarrival-gap/gap_keyform_incontainer.csv
\newcommand{\GapFiftyMsKeyRatio}{6.6}% data/runs/MECHANISM-2026-09-14-interarrival-gap/gap_keyform_incontainer.csv
\newcommand{\GapFiftyMsPreparsed}{157.01}% data/runs/MECHANISM-2026-09-14-interarrival-gap/gap_keyform_incontainer.csv
\newcommand{\GapFiftyMsString}{1037.51}% data/runs/MECHANISM-2026-09-14-interarrival-gap/gap_keyform_incontainer.csv
\newcommand{\GapTenMsKeyRatio}{6.7}% data/runs/MECHANISM-2026-09-14-interarrival-gap/gap_keyform_incontainer.csv
\newcommand{\GapTenMsPreparsed}{111.49}% data/runs/MECHANISM-2026-09-14-interarrival-gap/gap_keyform_incontainer.csv
\newcommand{\GapTenMsString}{747.82}% data/runs/MECHANISM-2026-09-14-interarrival-gap/gap_keyform_incontainer.csv
\newcommand{\GapTwoHundredMsKeyRatio}{4.8}% data/runs/MECHANISM-2026-09-14-interarrival-gap/gap_keyform_incontainer.csv
\newcommand{\GapTwoHundredMsPreparsed}{233.25}% data/runs/MECHANISM-2026-09-14-interarrival-gap/gap_keyform_incontainer.csv
\newcommand{\GapTwoHundredMsString}{1125.88}% data/runs/MECHANISM-2026-09-14-interarrival-gap/gap_keyform_incontainer.csv
\newcommand{\GapTwoMsKeyRatio}{7.3}% data/runs/MECHANISM-2026-09-14-interarrival-gap/gap_keyform_incontainer.csv
\newcommand{\GapTwoMsPreparsed}{87.73}% data/runs/MECHANISM-2026-09-14-interarrival-gap/gap_keyform_incontainer.csv
\newcommand{\GapTwoMsString}{637.32}% data/runs/MECHANISM-2026-09-14-interarrival-gap/gap_keyform_incontainer.csv
\newcommand{\Hundred}{100}% data/runs/2026-09-14T15-01-36Z-verifyrate-hs256-preparsed/run_metadata.json
\newcommand{\InsituAppSharePct}{28}% data/runs/INSITU-2026-09-14-verify-cost/openloop_ramp.csv
\newcommand{\InsituClientMean}{1944}% data/runs/INSITU-2026-09-14-verify-cost/openloop_ramp.csv
\newcommand{\InsituHandlerMean}{548.98}% data/runs/INSITU-2026-09-14-verify-cost/after.prom
\newcommand{\InsituVerifyMean}{421.21}% data/runs/INSITU-2026-09-14-verify-cost/after.prom
\newcommand{\InsituVersusMicrobenchRatio}{16.2}% data/runs/INSITU-2026-09-14-verify-cost/after.prom + data/runs/2026-09-13T20-34-41Z-microbench/microbench.csv
\newcommand{\KfAlgRatioPreparsed}{3.9}% data/runs/KEYFORM-2026-09-14/keyform_incontainer.csv
\newcommand{\KfAlgRatioString}{0.39}% data/runs/KEYFORM-2026-09-14/keyform_incontainer.csv
\newcommand{\KfHsConvPct}{98.1}% data/runs/KEYFORM-2026-09-14/keyform_incontainer.csv
\newcommand{\KfHsPreparsed}{8.26}% data/runs/KEYFORM-2026-09-14/keyform_incontainer.csv
\newcommand{\KfHsPreparsedIterations}{40000}% data/runs/KEYFORM-2026-09-14/keyform_incontainer.csv
\newcommand{\KfHsPreparsedPFifty}{7.17}% data/runs/KEYFORM-2026-09-14/keyform_incontainer.csv
\newcommand{\KfHsRatio}{52.8}% data/runs/KEYFORM-2026-09-14/keyform_incontainer.csv
\newcommand{\KfHsString}{435.72}% data/runs/KEYFORM-2026-09-14/keyform_incontainer.csv
\newcommand{\KfHsStringIterations}{40000}% data/runs/KEYFORM-2026-09-14/keyform_incontainer.csv
\newcommand{\KfHsStringPFifty}{421.17}% data/runs/KEYFORM-2026-09-14/keyform_incontainer.csv
\newcommand{\KfRsConvPct}{81.4}% data/runs/KEYFORM-2026-09-14/keyform_incontainer.csv
\newcommand{\KfRsPreparsed}{31.99}% data/runs/KEYFORM-2026-09-14/keyform_incontainer.csv
\newcommand{\KfRsPreparsedIterations}{40000}% data/runs/KEYFORM-2026-09-14/keyform_incontainer.csv
\newcommand{\KfRsPreparsedPFifty}{29.42}% data/runs/KEYFORM-2026-09-14/keyform_incontainer.csv
\newcommand{\KfRsRatio}{5.4}% data/runs/KEYFORM-2026-09-14/keyform_incontainer.csv
\newcommand{\KfRsString}{171.80}% data/runs/KEYFORM-2026-09-14/keyform_incontainer.csv
\newcommand{\KfRsStringIterations}{40000}% data/runs/KEYFORM-2026-09-14/keyform_incontainer.csv
\newcommand{\KfRsStringPFifty}{162.33}% data/runs/KEYFORM-2026-09-14/keyform_incontainer.csv
\newcommand{\MbFullVerifySpreadPct}{2.7}% data/runs/2026-09-13T20-34-41Z-microbench/microbench.spread.txt
\newcommand{\MbHsFourKbDecode}{2.090}% data/runs/2026-09-13T20-34-41Z-microbench/microbench.csv
\newcommand{\MbHsFourKbFull}{31.40}% data/runs/2026-09-13T20-34-41Z-microbench/microbench.csv
\newcommand{\MbHsFourKbIterations}{180000}% data/runs/2026-09-13T20-34-41Z-microbench/microbench.csv
\newcommand{\MbHsFourKbPrimitive}{2.613}% data/runs/2026-09-13T20-34-41Z-microbench/microbench.csv
\newcommand{\MbHsFourKbPrimitivePct}{8.3}% data/runs/2026-09-13T20-34-41Z-microbench/microbench.csv
\newcommand{\MbHsFourKbResidual}{26.70}% data/runs/2026-09-13T20-34-41Z-microbench/microbench.csv
\newcommand{\MbHsFourKbResidualPct}{85.0}% data/runs/2026-09-13T20-34-41Z-microbench/microbench.csv
\newcommand{\MbHsHalfKbDecode}{0.895}% data/runs/2026-09-13T20-34-41Z-microbench/microbench.csv
\newcommand{\MbHsHalfKbFull}{26.01}% data/runs/2026-09-13T20-34-41Z-microbench/microbench.csv
\newcommand{\MbHsHalfKbIterations}{180000}% data/runs/2026-09-13T20-34-41Z-microbench/microbench.csv
\newcommand{\MbHsHalfKbPrimitive}{1.484}% data/runs/2026-09-13T20-34-41Z-microbench/microbench.csv
\newcommand{\MbHsHalfKbPrimitivePct}{5.7}% data/runs/2026-09-13T20-34-41Z-microbench/microbench.csv
\newcommand{\MbHsHalfKbResidual}{23.63}% data/runs/2026-09-13T20-34-41Z-microbench/microbench.csv
\newcommand{\MbHsHalfKbResidualPct}{90.9}% data/runs/2026-09-13T20-34-41Z-microbench/microbench.csv
\newcommand{\MbHsTwoKbDecode}{1.355}% data/runs/2026-09-13T20-34-41Z-microbench/microbench.csv
\newcommand{\MbHsTwoKbFull}{28.67}% data/runs/2026-09-13T20-34-41Z-microbench/microbench.csv
\newcommand{\MbHsTwoKbIterations}{180000}% data/runs/2026-09-13T20-34-41Z-microbench/microbench.csv
\newcommand{\MbHsTwoKbPrimitive}{1.904}% data/runs/2026-09-13T20-34-41Z-microbench/microbench.csv
\newcommand{\MbHsTwoKbPrimitivePct}{6.6}% data/runs/2026-09-13T20-34-41Z-microbench/microbench.csv
\newcommand{\MbHsTwoKbResidual}{25.41}% data/runs/2026-09-13T20-34-41Z-microbench/microbench.csv
\newcommand{\MbHsTwoKbResidualPct}{88.6}% data/runs/2026-09-13T20-34-41Z-microbench/microbench.csv
\newcommand{\MbRsFourKbDecode}{2.009}% data/runs/2026-09-13T20-34-41Z-microbench/microbench.csv
\newcommand{\MbRsFourKbFull}{34.31}% data/runs/2026-09-13T20-34-41Z-microbench/microbench.csv
\newcommand{\MbRsFourKbIterations}{180000}% data/runs/2026-09-13T20-34-41Z-microbench/microbench.csv
\newcommand{\MbRsFourKbPrimitive}{12.422}% data/runs/2026-09-13T20-34-41Z-microbench/microbench.csv
\newcommand{\MbRsFourKbPrimitivePct}{36.2}% data/runs/2026-09-13T20-34-41Z-microbench/microbench.csv
\newcommand{\MbRsFourKbResidual}{19.88}% data/runs/2026-09-13T20-34-41Z-microbench/microbench.csv
\newcommand{\MbRsFourKbResidualPct}{57.9}% data/runs/2026-09-13T20-34-41Z-microbench/microbench.csv
\newcommand{\MbRsHalfKbDecode}{0.894}% data/runs/2026-09-13T20-34-41Z-microbench/microbench.csv
\newcommand{\MbRsHalfKbFull}{29.00}% data/runs/2026-09-13T20-34-41Z-microbench/microbench.csv
\newcommand{\MbRsHalfKbIterations}{180000}% data/runs/2026-09-13T20-34-41Z-microbench/microbench.csv
\newcommand{\MbRsHalfKbPrimitive}{11.159}% data/runs/2026-09-13T20-34-41Z-microbench/microbench.csv
\newcommand{\MbRsHalfKbPrimitivePct}{38.5}% data/runs/2026-09-13T20-34-41Z-microbench/microbench.csv
\newcommand{\MbRsHalfKbResidual}{16.95}% data/runs/2026-09-13T20-34-41Z-microbench/microbench.csv
\newcommand{\MbRsHalfKbResidualPct}{58.4}% data/runs/2026-09-13T20-34-41Z-microbench/microbench.csv
\newcommand{\MbRsTwoKbDecode}{1.353}% data/runs/2026-09-13T20-34-41Z-microbench/microbench.csv
\newcommand{\MbRsTwoKbFull}{30.90}% data/runs/2026-09-13T20-34-41Z-microbench/microbench.csv
\newcommand{\MbRsTwoKbIterations}{180000}% data/runs/2026-09-13T20-34-41Z-microbench/microbench.csv
\newcommand{\MbRsTwoKbPrimitive}{12.028}% data/runs/2026-09-13T20-34-41Z-microbench/microbench.csv
\newcommand{\MbRsTwoKbPrimitivePct}{38.9}% data/runs/2026-09-13T20-34-41Z-microbench/microbench.csv
\newcommand{\MbRsTwoKbResidual}{17.51}% data/runs/2026-09-13T20-34-41Z-microbench/microbench.csv
\newcommand{\MbRsTwoKbResidualPct}{56.7}% data/runs/2026-09-13T20-34-41Z-microbench/microbench.csv
\newcommand{\MbWorstStageSpreadPct}{12.6}% data/runs/2026-09-13T20-34-41Z-microbench/microbench.spread.txt
\newcommand{\Ten}{10}% data/runs/2026-09-14T15-01-36Z-verifyrate-hs256-preparsed/run_metadata.json
\newcommand{\TightKeyRatio}{43.1}% data/runs/MECHANISM-2026-09-14-interarrival-gap/gap_keyform_incontainer.csv
\newcommand{\TightPreparsed}{9.63}% data/runs/MECHANISM-2026-09-14-interarrival-gap/gap_keyform_incontainer.csv
\newcommand{\TightString}{415.43}% data/runs/MECHANISM-2026-09-14-interarrival-gap/gap_keyform_incontainer.csv
\newcommand{\TwentyFive}{25}% data/runs/2026-09-14T15-01-36Z-verifyrate-hs256-preparsed/run_metadata.json
\newcommand{\TwoHundred}{200}% data/runs/2026-09-14T15-01-36Z-verifyrate-hs256-preparsed/run_metadata.json
\newcommand{\VrAlgRatioFifty}{1.60}% data/runs/2026-09-14T06-41-56Z-verifyrate-rs256/verify_rate_curve.csv
\newcommand{\VrAlgRatioFive}{1.59}% data/runs/2026-09-14T06-41-56Z-verifyrate-rs256/verify_rate_curve.csv
\newcommand{\VrAlgRatioFourHundred}{1.50}% data/runs/2026-09-14T06-41-56Z-verifyrate-rs256/verify_rate_curve.csv
\newcommand{\VrAlgRatioHundred}{1.51}% data/runs/2026-09-14T06-41-56Z-verifyrate-rs256/verify_rate_curve.csv
\newcommand{\VrAlgRatioTen}{1.44}% data/runs/2026-09-14T06-41-56Z-verifyrate-rs256/verify_rate_curve.csv
\newcommand{\VrAlgRatioTwentyFive}{1.61}% data/runs/2026-09-14T06-41-56Z-verifyrate-rs256/verify_rate_curve.csv
\newcommand{\VrAlgRatioTwoHundred}{1.03}% data/runs/2026-09-14T06-41-56Z-verifyrate-rs256/verify_rate_curve.csv
\newcommand{\VrConvPctFifty}{83}% data/runs/2026-09-14T05-57-22Z-verifyrate-hs256/verify_rate_curve.csv
\newcommand{\VrConvPctFive}{76}% data/runs/2026-09-14T05-57-22Z-verifyrate-hs256/verify_rate_curve.csv
\newcommand{\VrConvPctFourHundred}{91}% data/runs/2026-09-14T05-57-22Z-verifyrate-hs256/verify_rate_curve.csv
\newcommand{\VrConvPctHundred}{86}% data/runs/2026-09-14T05-57-22Z-verifyrate-hs256/verify_rate_curve.csv
\newcommand{\VrConvPctTen}{76}% data/runs/2026-09-14T05-57-22Z-verifyrate-hs256/verify_rate_curve.csv
\newcommand{\VrConvPctTwentyFive}{82}% data/runs/2026-09-14T05-57-22Z-verifyrate-hs256/verify_rate_curve.csv
\newcommand{\VrConvPctTwoHundred}{86}% data/runs/2026-09-14T05-57-22Z-verifyrate-hs256/verify_rate_curve.csv
\newcommand{\VrHsPreparsedAscFiftyCount}{9001}% data/runs/2026-09-14T15-01-36Z-verifyrate-hs256-preparsed/verify_rate_curve.csv
\newcommand{\VrHsPreparsedAscFiftyLoad}{2.70}% data/runs/2026-09-14T15-01-36Z-verifyrate-hs256-preparsed/verify_rate_curve.csv
\newcommand{\VrHsPreparsedAscFiftyMean}{127.19}% data/runs/2026-09-14T15-01-36Z-verifyrate-hs256-preparsed/verify_rate_curve.csv
\newcommand{\VrHsPreparsedAscFiftyPFifty}{100.26}% data/runs/2026-09-14T15-01-36Z-verifyrate-hs256-preparsed/verify_rate_curve.csv
\newcommand{\VrHsPreparsedAscFiftyPNinetyNine}{482.59}% data/runs/2026-09-14T15-01-36Z-verifyrate-hs256-preparsed/verify_rate_curve.csv
\newcommand{\VrHsPreparsedAscFiveCount}{901}% data/runs/2026-09-14T15-01-36Z-verifyrate-hs256-preparsed/verify_rate_curve.csv
\newcommand{\VrHsPreparsedAscFiveLoad}{3.95}% data/runs/2026-09-14T15-01-36Z-verifyrate-hs256-preparsed/verify_rate_curve.csv
\newcommand{\VrHsPreparsedAscFiveMean}{221.42}% data/runs/2026-09-14T15-01-36Z-verifyrate-hs256-preparsed/verify_rate_curve.csv
\newcommand{\VrHsPreparsedAscFivePFifty}{198.83}% data/runs/2026-09-14T15-01-36Z-verifyrate-hs256-preparsed/verify_rate_curve.csv
\newcommand{\VrHsPreparsedAscFivePNinetyNine}{597.25}% data/runs/2026-09-14T15-01-36Z-verifyrate-hs256-preparsed/verify_rate_curve.csv
\newcommand{\VrHsPreparsedAscFourHundredCount}{72001}% data/runs/2026-09-14T15-01-36Z-verifyrate-hs256-preparsed/verify_rate_curve.csv
\newcommand{\VrHsPreparsedAscFourHundredLoad}{5.33}% data/runs/2026-09-14T15-01-36Z-verifyrate-hs256-preparsed/verify_rate_curve.csv
\newcommand{\VrHsPreparsedAscFourHundredMean}{39.38}% data/runs/2026-09-14T15-01-36Z-verifyrate-hs256-preparsed/verify_rate_curve.csv
\newcommand{\VrHsPreparsedAscFourHundredPFifty}{32.18}% data/runs/2026-09-14T15-01-36Z-verifyrate-hs256-preparsed/verify_rate_curve.csv
\newcommand{\VrHsPreparsedAscFourHundredPNinetyNine}{243.73}% data/runs/2026-09-14T15-01-36Z-verifyrate-hs256-preparsed/verify_rate_curve.csv
\newcommand{\VrHsPreparsedAscHundredCount}{17993}% data/runs/2026-09-14T15-01-36Z-verifyrate-hs256-preparsed/verify_rate_curve.csv
\newcommand{\VrHsPreparsedAscHundredLoad}{2.41}% data/runs/2026-09-14T15-01-36Z-verifyrate-hs256-preparsed/verify_rate_curve.csv
\newcommand{\VrHsPreparsedAscHundredMean}{79.70}% data/runs/2026-09-14T15-01-36Z-verifyrate-hs256-preparsed/verify_rate_curve.csv
\newcommand{\VrHsPreparsedAscHundredPFifty}{59.05}% data/runs/2026-09-14T15-01-36Z-verifyrate-hs256-preparsed/verify_rate_curve.csv
\newcommand{\VrHsPreparsedAscHundredPNinetyNine}{358.02}% data/runs/2026-09-14T15-01-36Z-verifyrate-hs256-preparsed/verify_rate_curve.csv
\newcommand{\VrHsPreparsedAscTenCount}{1801}% data/runs/2026-09-14T15-01-36Z-verifyrate-hs256-preparsed/verify_rate_curve.csv
\newcommand{\VrHsPreparsedAscTenLoad}{2.58}% data/runs/2026-09-14T15-01-36Z-verifyrate-hs256-preparsed/verify_rate_curve.csv
\newcommand{\VrHsPreparsedAscTenMean}{235.01}% data/runs/2026-09-14T15-01-36Z-verifyrate-hs256-preparsed/verify_rate_curve.csv
\newcommand{\VrHsPreparsedAscTenPFifty}{205.01}% data/runs/2026-09-14T15-01-36Z-verifyrate-hs256-preparsed/verify_rate_curve.csv
\newcommand{\VrHsPreparsedAscTenPNinetyNine}{1058.47}% data/runs/2026-09-14T15-01-36Z-verifyrate-hs256-preparsed/verify_rate_curve.csv
\newcommand{\VrHsPreparsedAscTwentyFiveCount}{4500}% data/runs/2026-09-14T15-01-36Z-verifyrate-hs256-preparsed/verify_rate_curve.csv
\newcommand{\VrHsPreparsedAscTwentyFiveLoad}{3.98}% data/runs/2026-09-14T15-01-36Z-verifyrate-hs256-preparsed/verify_rate_curve.csv
\newcommand{\VrHsPreparsedAscTwentyFiveMean}{168.44}% data/runs/2026-09-14T15-01-36Z-verifyrate-hs256-preparsed/verify_rate_curve.csv
\newcommand{\VrHsPreparsedAscTwentyFivePFifty}{145.68}% data/runs/2026-09-14T15-01-36Z-verifyrate-hs256-preparsed/verify_rate_curve.csv
\newcommand{\VrHsPreparsedAscTwentyFivePNinetyNine}{569.10}% data/runs/2026-09-14T15-01-36Z-verifyrate-hs256-preparsed/verify_rate_curve.csv
\newcommand{\VrHsPreparsedAscTwoHundredCount}{36001}% data/runs/2026-09-14T15-01-36Z-verifyrate-hs256-preparsed/verify_rate_curve.csv
\newcommand{\VrHsPreparsedAscTwoHundredLoad}{4.88}% data/runs/2026-09-14T15-01-36Z-verifyrate-hs256-preparsed/verify_rate_curve.csv
\newcommand{\VrHsPreparsedAscTwoHundredMean}{71.68}% data/runs/2026-09-14T15-01-36Z-verifyrate-hs256-preparsed/verify_rate_curve.csv
\newcommand{\VrHsPreparsedAscTwoHundredPFifty}{49.63}% data/runs/2026-09-14T15-01-36Z-verifyrate-hs256-preparsed/verify_rate_curve.csv
\newcommand{\VrHsPreparsedAscTwoHundredPNinetyNine}{408.52}% data/runs/2026-09-14T15-01-36Z-verifyrate-hs256-preparsed/verify_rate_curve.csv
\newcommand{\VrHsPreparsedDescFiftyCount}{9001}% data/runs/2026-09-14T15-01-36Z-verifyrate-hs256-preparsed/verify_rate_curve.csv
\newcommand{\VrHsPreparsedDescFiftyLoad}{2.53}% data/runs/2026-09-14T15-01-36Z-verifyrate-hs256-preparsed/verify_rate_curve.csv
\newcommand{\VrHsPreparsedDescFiftyMean}{116.67}% data/runs/2026-09-14T15-01-36Z-verifyrate-hs256-preparsed/verify_rate_curve.csv
\newcommand{\VrHsPreparsedDescFiftyPFifty}{88.18}% data/runs/2026-09-14T15-01-36Z-verifyrate-hs256-preparsed/verify_rate_curve.csv
\newcommand{\VrHsPreparsedDescFiftyPNinetyNine}{327.55}% data/runs/2026-09-14T15-01-36Z-verifyrate-hs256-preparsed/verify_rate_curve.csv
\newcommand{\VrHsPreparsedDescFiveCount}{901}% data/runs/2026-09-14T15-01-36Z-verifyrate-hs256-preparsed/verify_rate_curve.csv
\newcommand{\VrHsPreparsedDescFiveLoad}{2.35}% data/runs/2026-09-14T15-01-36Z-verifyrate-hs256-preparsed/verify_rate_curve.csv
\newcommand{\VrHsPreparsedDescFiveMean}{200.88}% data/runs/2026-09-14T15-01-36Z-verifyrate-hs256-preparsed/verify_rate_curve.csv
\newcommand{\VrHsPreparsedDescFivePFifty}{179.73}% data/runs/2026-09-14T15-01-36Z-verifyrate-hs256-preparsed/verify_rate_curve.csv
\newcommand{\VrHsPreparsedDescFivePNinetyNine}{899.57}% data/runs/2026-09-14T15-01-36Z-verifyrate-hs256-preparsed/verify_rate_curve.csv
\newcommand{\VrHsPreparsedDescFourHundredCount}{71987}% data/runs/2026-09-14T15-01-36Z-verifyrate-hs256-preparsed/verify_rate_curve.csv
\newcommand{\VrHsPreparsedDescFourHundredLoad}{4.21}% data/runs/2026-09-14T15-01-36Z-verifyrate-hs256-preparsed/verify_rate_curve.csv
\newcommand{\VrHsPreparsedDescFourHundredMean}{34.61}% data/runs/2026-09-14T15-01-36Z-verifyrate-hs256-preparsed/verify_rate_curve.csv
\newcommand{\VrHsPreparsedDescFourHundredPFifty}{30.93}% data/runs/2026-09-14T15-01-36Z-verifyrate-hs256-preparsed/verify_rate_curve.csv
\newcommand{\VrHsPreparsedDescFourHundredPNinetyNine}{147.41}% data/runs/2026-09-14T15-01-36Z-verifyrate-hs256-preparsed/verify_rate_curve.csv
\newcommand{\VrHsPreparsedDescHundredCount}{18001}% data/runs/2026-09-14T15-01-36Z-verifyrate-hs256-preparsed/verify_rate_curve.csv
\newcommand{\VrHsPreparsedDescHundredLoad}{3.38}% data/runs/2026-09-14T15-01-36Z-verifyrate-hs256-preparsed/verify_rate_curve.csv
\newcommand{\VrHsPreparsedDescHundredMean}{66.93}% data/runs/2026-09-14T15-01-36Z-verifyrate-hs256-preparsed/verify_rate_curve.csv
\newcommand{\VrHsPreparsedDescHundredPFifty}{53.65}% data/runs/2026-09-14T15-01-36Z-verifyrate-hs256-preparsed/verify_rate_curve.csv
\newcommand{\VrHsPreparsedDescHundredPNinetyNine}{286.62}% data/runs/2026-09-14T15-01-36Z-verifyrate-hs256-preparsed/verify_rate_curve.csv
\newcommand{\VrHsPreparsedDescTenCount}{1801}% data/runs/2026-09-14T15-01-36Z-verifyrate-hs256-preparsed/verify_rate_curve.csv
\newcommand{\VrHsPreparsedDescTenLoad}{2.36}% data/runs/2026-09-14T15-01-36Z-verifyrate-hs256-preparsed/verify_rate_curve.csv
\newcommand{\VrHsPreparsedDescTenMean}{188.06}% data/runs/2026-09-14T15-01-36Z-verifyrate-hs256-preparsed/verify_rate_curve.csv
\newcommand{\VrHsPreparsedDescTenPFifty}{169.50}% data/runs/2026-09-14T15-01-36Z-verifyrate-hs256-preparsed/verify_rate_curve.csv
\newcommand{\VrHsPreparsedDescTenPNinetyNine}{879.60}% data/runs/2026-09-14T15-01-36Z-verifyrate-hs256-preparsed/verify_rate_curve.csv
\newcommand{\VrHsPreparsedDescTwentyFiveCount}{4501}% data/runs/2026-09-14T15-01-36Z-verifyrate-hs256-preparsed/verify_rate_curve.csv
\newcommand{\VrHsPreparsedDescTwentyFiveLoad}{1.95}% data/runs/2026-09-14T15-01-36Z-verifyrate-hs256-preparsed/verify_rate_curve.csv
\newcommand{\VrHsPreparsedDescTwentyFiveMean}{160.58}% data/runs/2026-09-14T15-01-36Z-verifyrate-hs256-preparsed/verify_rate_curve.csv
\newcommand{\VrHsPreparsedDescTwentyFivePFifty}{137.47}% data/runs/2026-09-14T15-01-36Z-verifyrate-hs256-preparsed/verify_rate_curve.csv
\newcommand{\VrHsPreparsedDescTwentyFivePNinetyNine}{557.76}% data/runs/2026-09-14T15-01-36Z-verifyrate-hs256-preparsed/verify_rate_curve.csv
\newcommand{\VrHsPreparsedDescTwoHundredCount}{36001}% data/runs/2026-09-14T15-01-36Z-verifyrate-hs256-preparsed/verify_rate_curve.csv
\newcommand{\VrHsPreparsedDescTwoHundredLoad}{4.90}% data/runs/2026-09-14T15-01-36Z-verifyrate-hs256-preparsed/verify_rate_curve.csv
\newcommand{\VrHsPreparsedDescTwoHundredMean}{42.78}% data/runs/2026-09-14T15-01-36Z-verifyrate-hs256-preparsed/verify_rate_curve.csv
\newcommand{\VrHsPreparsedDescTwoHundredPFifty}{34.30}% data/runs/2026-09-14T15-01-36Z-verifyrate-hs256-preparsed/verify_rate_curve.csv
\newcommand{\VrHsPreparsedDescTwoHundredPNinetyNine}{239.34}% data/runs/2026-09-14T15-01-36Z-verifyrate-hs256-preparsed/verify_rate_curve.csv
\newcommand{\VrHsStringAscFiftyCount}{9000}% data/runs/2026-09-14T05-57-22Z-verifyrate-hs256/verify_rate_curve.csv
\newcommand{\VrHsStringAscFiftyLoad}{1.93}% data/runs/2026-09-14T05-57-22Z-verifyrate-hs256/verify_rate_curve.csv
\newcommand{\VrHsStringAscFiftyMean}{758.57}% data/runs/2026-09-14T05-57-22Z-verifyrate-hs256/verify_rate_curve.csv
\newcommand{\VrHsStringAscFiftyPFifty}{713.77}% data/runs/2026-09-14T05-57-22Z-verifyrate-hs256/verify_rate_curve.csv
\newcommand{\VrHsStringAscFiftyPNinetyNine}{2236.45}% data/runs/2026-09-14T05-57-22Z-verifyrate-hs256/verify_rate_curve.csv
\newcommand{\VrHsStringAscFiveCount}{901}% data/runs/2026-09-14T05-57-22Z-verifyrate-hs256/verify_rate_curve.csv
\newcommand{\VrHsStringAscFiveLoad}{5.47}% data/runs/2026-09-14T05-57-22Z-verifyrate-hs256/verify_rate_curve.csv
\newcommand{\VrHsStringAscFiveMean}{928.24}% data/runs/2026-09-14T05-57-22Z-verifyrate-hs256/verify_rate_curve.csv
\newcommand{\VrHsStringAscFivePFifty}{955.20}% data/runs/2026-09-14T05-57-22Z-verifyrate-hs256/verify_rate_curve.csv
\newcommand{\VrHsStringAscFivePNinetyNine}{2449.29}% data/runs/2026-09-14T05-57-22Z-verifyrate-hs256/verify_rate_curve.csv
\newcommand{\VrHsStringAscFourHundredCount}{72001}% data/runs/2026-09-14T05-57-22Z-verifyrate-hs256/verify_rate_curve.csv
\newcommand{\VrHsStringAscFourHundredLoad}{2.95}% data/runs/2026-09-14T05-57-22Z-verifyrate-hs256/verify_rate_curve.csv
\newcommand{\VrHsStringAscFourHundredMean}{452.82}% data/runs/2026-09-14T05-57-22Z-verifyrate-hs256/verify_rate_curve.csv
\newcommand{\VrHsStringAscFourHundredPFifty}{451.95}% data/runs/2026-09-14T05-57-22Z-verifyrate-hs256/verify_rate_curve.csv
\newcommand{\VrHsStringAscFourHundredPNinetyNine}{737.79}% data/runs/2026-09-14T05-57-22Z-verifyrate-hs256/verify_rate_curve.csv
\newcommand{\VrHsStringAscHundredCount}{18001}% data/runs/2026-09-14T05-57-22Z-verifyrate-hs256/verify_rate_curve.csv
\newcommand{\VrHsStringAscHundredLoad}{1.96}% data/runs/2026-09-14T05-57-22Z-verifyrate-hs256/verify_rate_curve.csv
\newcommand{\VrHsStringAscHundredMean}{571.59}% data/runs/2026-09-14T05-57-22Z-verifyrate-hs256/verify_rate_curve.csv
\newcommand{\VrHsStringAscHundredPFifty}{537.71}% data/runs/2026-09-14T05-57-22Z-verifyrate-hs256/verify_rate_curve.csv
\newcommand{\VrHsStringAscHundredPNinetyNine}{1184.80}% data/runs/2026-09-14T05-57-22Z-verifyrate-hs256/verify_rate_curve.csv
\newcommand{\VrHsStringAscTenCount}{1801}% data/runs/2026-09-14T05-57-22Z-verifyrate-hs256/verify_rate_curve.csv
\newcommand{\VrHsStringAscTenLoad}{2.77}% data/runs/2026-09-14T05-57-22Z-verifyrate-hs256/verify_rate_curve.csv
\newcommand{\VrHsStringAscTenMean}{976.28}% data/runs/2026-09-14T05-57-22Z-verifyrate-hs256/verify_rate_curve.csv
\newcommand{\VrHsStringAscTenPFifty}{1022.66}% data/runs/2026-09-14T05-57-22Z-verifyrate-hs256/verify_rate_curve.csv
\newcommand{\VrHsStringAscTenPNinetyNine}{2464.28}% data/runs/2026-09-14T05-57-22Z-verifyrate-hs256/verify_rate_curve.csv
\newcommand{\VrHsStringAscTwentyFiveCount}{4501}% data/runs/2026-09-14T05-57-22Z-verifyrate-hs256/verify_rate_curve.csv
\newcommand{\VrHsStringAscTwentyFiveLoad}{3.05}% data/runs/2026-09-14T05-57-22Z-verifyrate-hs256/verify_rate_curve.csv
\newcommand{\VrHsStringAscTwentyFiveMean}{918.54}% data/runs/2026-09-14T05-57-22Z-verifyrate-hs256/verify_rate_curve.csv
\newcommand{\VrHsStringAscTwentyFivePFifty}{906.76}% data/runs/2026-09-14T05-57-22Z-verifyrate-hs256/verify_rate_curve.csv
\newcommand{\VrHsStringAscTwentyFivePNinetyNine}{2461.02}% data/runs/2026-09-14T05-57-22Z-verifyrate-hs256/verify_rate_curve.csv
\newcommand{\VrHsStringAscTwoHundredCount}{36000}% data/runs/2026-09-14T05-57-22Z-verifyrate-hs256/verify_rate_curve.csv
\newcommand{\VrHsStringAscTwoHundredLoad}{2.05}% data/runs/2026-09-14T05-57-22Z-verifyrate-hs256/verify_rate_curve.csv
\newcommand{\VrHsStringAscTwoHundredMean}{494.91}% data/runs/2026-09-14T05-57-22Z-verifyrate-hs256/verify_rate_curve.csv
\newcommand{\VrHsStringAscTwoHundredPFifty}{470.82}% data/runs/2026-09-14T05-57-22Z-verifyrate-hs256/verify_rate_curve.csv
\newcommand{\VrHsStringAscTwoHundredPNinetyNine}{1153.04}% data/runs/2026-09-14T05-57-22Z-verifyrate-hs256/verify_rate_curve.csv
\newcommand{\VrHsStringDescFiftyCount}{9001}% data/runs/2026-09-14T05-57-22Z-verifyrate-hs256/verify_rate_curve.csv
\newcommand{\VrHsStringDescFiftyLoad}{2.31}% data/runs/2026-09-14T05-57-22Z-verifyrate-hs256/verify_rate_curve.csv
\newcommand{\VrHsStringDescFiftyMean}{737.12}% data/runs/2026-09-14T05-57-22Z-verifyrate-hs256/verify_rate_curve.csv
\newcommand{\VrHsStringDescFiftyPFifty}{758.26}% data/runs/2026-09-14T05-57-22Z-verifyrate-hs256/verify_rate_curve.csv
\newcommand{\VrHsStringDescFiftyPNinetyNine}{1742.30}% data/runs/2026-09-14T05-57-22Z-verifyrate-hs256/verify_rate_curve.csv
\newcommand{\VrHsStringDescFiveCount}{901}% data/runs/2026-09-14T05-57-22Z-verifyrate-hs256/verify_rate_curve.csv
\newcommand{\VrHsStringDescFiveLoad}{4.47}% data/runs/2026-09-14T05-57-22Z-verifyrate-hs256/verify_rate_curve.csv
\newcommand{\VrHsStringDescFiveMean}{1027.97}% data/runs/2026-09-14T05-57-22Z-verifyrate-hs256/verify_rate_curve.csv
\newcommand{\VrHsStringDescFivePFifty}{1120.35}% data/runs/2026-09-14T05-57-22Z-verifyrate-hs256/verify_rate_curve.csv
\newcommand{\VrHsStringDescFivePNinetyNine}{2469.34}% data/runs/2026-09-14T05-57-22Z-verifyrate-hs256/verify_rate_curve.csv
\newcommand{\VrHsStringDescFourHundredCount}{71929}% data/runs/2026-09-14T05-57-22Z-verifyrate-hs256/verify_rate_curve.csv
\newcommand{\VrHsStringDescFourHundredLoad}{2.93}% data/runs/2026-09-14T05-57-22Z-verifyrate-hs256/verify_rate_curve.csv
\newcommand{\VrHsStringDescFourHundredMean}{467.62}% data/runs/2026-09-14T05-57-22Z-verifyrate-hs256/verify_rate_curve.csv
\newcommand{\VrHsStringDescFourHundredPFifty}{453.31}% data/runs/2026-09-14T05-57-22Z-verifyrate-hs256/verify_rate_curve.csv
\newcommand{\VrHsStringDescFourHundredPNinetyNine}{970.54}% data/runs/2026-09-14T05-57-22Z-verifyrate-hs256/verify_rate_curve.csv
\newcommand{\VrHsStringDescHundredCount}{18001}% data/runs/2026-09-14T05-57-22Z-verifyrate-hs256/verify_rate_curve.csv
\newcommand{\VrHsStringDescHundredLoad}{2.63}% data/runs/2026-09-14T05-57-22Z-verifyrate-hs256/verify_rate_curve.csv
\newcommand{\VrHsStringDescHundredMean}{607.78}% data/runs/2026-09-14T05-57-22Z-verifyrate-hs256/verify_rate_curve.csv
\newcommand{\VrHsStringDescHundredPFifty}{562.23}% data/runs/2026-09-14T05-57-22Z-verifyrate-hs256/verify_rate_curve.csv
\newcommand{\VrHsStringDescHundredPNinetyNine}{1188.73}% data/runs/2026-09-14T05-57-22Z-verifyrate-hs256/verify_rate_curve.csv
\newcommand{\VrHsStringDescTenCount}{1800}% data/runs/2026-09-14T05-57-22Z-verifyrate-hs256/verify_rate_curve.csv
\newcommand{\VrHsStringDescTenLoad}{2.42}% data/runs/2026-09-14T05-57-22Z-verifyrate-hs256/verify_rate_curve.csv
\newcommand{\VrHsStringDescTenMean}{1017.34}% data/runs/2026-09-14T05-57-22Z-verifyrate-hs256/verify_rate_curve.csv
\newcommand{\VrHsStringDescTenPFifty}{1047.25}% data/runs/2026-09-14T05-57-22Z-verifyrate-hs256/verify_rate_curve.csv
\newcommand{\VrHsStringDescTenPNinetyNine}{2495.71}% data/runs/2026-09-14T05-57-22Z-verifyrate-hs256/verify_rate_curve.csv
\newcommand{\VrHsStringDescTwentyFiveCount}{4500}% data/runs/2026-09-14T05-57-22Z-verifyrate-hs256/verify_rate_curve.csv
\newcommand{\VrHsStringDescTwentyFiveLoad}{2.60}% data/runs/2026-09-14T05-57-22Z-verifyrate-hs256/verify_rate_curve.csv
\newcommand{\VrHsStringDescTwentyFiveMean}{934.47}% data/runs/2026-09-14T05-57-22Z-verifyrate-hs256/verify_rate_curve.csv
\newcommand{\VrHsStringDescTwentyFivePFifty}{988.97}% data/runs/2026-09-14T05-57-22Z-verifyrate-hs256/verify_rate_curve.csv
\newcommand{\VrHsStringDescTwentyFivePNinetyNine}{2461.91}% data/runs/2026-09-14T05-57-22Z-verifyrate-hs256/verify_rate_curve.csv
\newcommand{\VrHsStringDescTwoHundredCount}{36000}% data/runs/2026-09-14T05-57-22Z-verifyrate-hs256/verify_rate_curve.csv
\newcommand{\VrHsStringDescTwoHundredLoad}{3.95}% data/runs/2026-09-14T05-57-22Z-verifyrate-hs256/verify_rate_curve.csv
\newcommand{\VrHsStringDescTwoHundredMean}{532.50}% data/runs/2026-09-14T05-57-22Z-verifyrate-hs256/verify_rate_curve.csv
\newcommand{\VrHsStringDescTwoHundredPFifty}{475.25}% data/runs/2026-09-14T05-57-22Z-verifyrate-hs256/verify_rate_curve.csv
\newcommand{\VrHsStringDescTwoHundredPNinetyNine}{1162.25}% data/runs/2026-09-14T05-57-22Z-verifyrate-hs256/verify_rate_curve.csv
\newcommand{\VrPassAgreementPct}{10.7}% data/runs/2026-09-14T05-57-22Z-verifyrate-hs256/verify_rate_curve.csv
\newcommand{\VrRatioFifty}{6.0}% data/runs/2026-09-14T05-57-22Z-verifyrate-hs256/verify_rate_curve.csv
\newcommand{\VrRatioFive}{4.2}% data/runs/2026-09-14T05-57-22Z-verifyrate-hs256/verify_rate_curve.csv
\newcommand{\VrRatioFourHundred}{11.5}% data/runs/2026-09-14T05-57-22Z-verifyrate-hs256/verify_rate_curve.csv
\newcommand{\VrRatioHundred}{7.2}% data/runs/2026-09-14T05-57-22Z-verifyrate-hs256/verify_rate_curve.csv
\newcommand{\VrRatioTen}{4.2}% data/runs/2026-09-14T05-57-22Z-verifyrate-hs256/verify_rate_curve.csv
\newcommand{\VrRatioTwentyFive}{5.5}% data/runs/2026-09-14T05-57-22Z-verifyrate-hs256/verify_rate_curve.csv
\newcommand{\VrRatioTwoHundred}{6.9}% data/runs/2026-09-14T05-57-22Z-verifyrate-hs256/verify_rate_curve.csv
\newcommand{\VrRsPreparsedAscFiftyCount}{9000}% data/runs/2026-09-14T06-41-56Z-verifyrate-rs256/verify_rate_curve.csv
\newcommand{\VrRsPreparsedAscFiftyLoad}{1.22}% data/runs/2026-09-14T06-41-56Z-verifyrate-rs256/verify_rate_curve.csv
\newcommand{\VrRsPreparsedAscFiftyMean}{203.19}% data/runs/2026-09-14T06-41-56Z-verifyrate-rs256/verify_rate_curve.csv
\newcommand{\VrRsPreparsedAscFiftyPFifty}{146.25}% data/runs/2026-09-14T06-41-56Z-verifyrate-rs256/verify_rate_curve.csv
\newcommand{\VrRsPreparsedAscFiftyPNinetyNine}{592.18}% data/runs/2026-09-14T06-41-56Z-verifyrate-rs256/verify_rate_curve.csv
\newcommand{\VrRsPreparsedAscFiveCount}{901}% data/runs/2026-09-14T06-41-56Z-verifyrate-rs256/verify_rate_curve.csv
\newcommand{\VrRsPreparsedAscFiveLoad}{2.77}% data/runs/2026-09-14T06-41-56Z-verifyrate-rs256/verify_rate_curve.csv
\newcommand{\VrRsPreparsedAscFiveMean}{351.05}% data/runs/2026-09-14T06-41-56Z-verifyrate-rs256/verify_rate_curve.csv
\newcommand{\VrRsPreparsedAscFivePFifty}{333.92}% data/runs/2026-09-14T06-41-56Z-verifyrate-rs256/verify_rate_curve.csv
\newcommand{\VrRsPreparsedAscFivePNinetyNine}{1066.33}% data/runs/2026-09-14T06-41-56Z-verifyrate-rs256/verify_rate_curve.csv
\newcommand{\VrRsPreparsedAscFourHundredCount}{72001}% data/runs/2026-09-14T06-41-56Z-verifyrate-rs256/verify_rate_curve.csv
\newcommand{\VrRsPreparsedAscFourHundredLoad}{4.02}% data/runs/2026-09-14T06-41-56Z-verifyrate-rs256/verify_rate_curve.csv
\newcommand{\VrRsPreparsedAscFourHundredMean}{59.11}% data/runs/2026-09-14T06-41-56Z-verifyrate-rs256/verify_rate_curve.csv
\newcommand{\VrRsPreparsedAscFourHundredPFifty}{61.35}% data/runs/2026-09-14T06-41-56Z-verifyrate-rs256/verify_rate_curve.csv
\newcommand{\VrRsPreparsedAscFourHundredPNinetyNine}{149.57}% data/runs/2026-09-14T06-41-56Z-verifyrate-rs256/verify_rate_curve.csv
\newcommand{\VrRsPreparsedAscHundredCount}{17989}% data/runs/2026-09-14T06-41-56Z-verifyrate-rs256/verify_rate_curve.csv
\newcommand{\VrRsPreparsedAscHundredLoad}{1.48}% data/runs/2026-09-14T06-41-56Z-verifyrate-rs256/verify_rate_curve.csv
\newcommand{\VrRsPreparsedAscHundredMean}{120.42}% data/runs/2026-09-14T06-41-56Z-verifyrate-rs256/verify_rate_curve.csv
\newcommand{\VrRsPreparsedAscHundredPFifty}{91.67}% data/runs/2026-09-14T06-41-56Z-verifyrate-rs256/verify_rate_curve.csv
\newcommand{\VrRsPreparsedAscHundredPNinetyNine}{551.55}% data/runs/2026-09-14T06-41-56Z-verifyrate-rs256/verify_rate_curve.csv
\newcommand{\VrRsPreparsedAscTenCount}{1801}% data/runs/2026-09-14T06-41-56Z-verifyrate-rs256/verify_rate_curve.csv
\newcommand{\VrRsPreparsedAscTenLoad}{2.74}% data/runs/2026-09-14T06-41-56Z-verifyrate-rs256/verify_rate_curve.csv
\newcommand{\VrRsPreparsedAscTenMean}{338.89}% data/runs/2026-09-14T06-41-56Z-verifyrate-rs256/verify_rate_curve.csv
\newcommand{\VrRsPreparsedAscTenPFifty}{321.54}% data/runs/2026-09-14T06-41-56Z-verifyrate-rs256/verify_rate_curve.csv
\newcommand{\VrRsPreparsedAscTenPNinetyNine}{1092.64}% data/runs/2026-09-14T06-41-56Z-verifyrate-rs256/verify_rate_curve.csv
\newcommand{\VrRsPreparsedAscTwentyFiveCount}{4500}% data/runs/2026-09-14T06-41-56Z-verifyrate-rs256/verify_rate_curve.csv
\newcommand{\VrRsPreparsedAscTwentyFiveLoad}{2.77}% data/runs/2026-09-14T06-41-56Z-verifyrate-rs256/verify_rate_curve.csv
\newcommand{\VrRsPreparsedAscTwentyFiveMean}{270.77}% data/runs/2026-09-14T06-41-56Z-verifyrate-rs256/verify_rate_curve.csv
\newcommand{\VrRsPreparsedAscTwentyFivePFifty}{278.22}% data/runs/2026-09-14T06-41-56Z-verifyrate-rs256/verify_rate_curve.csv
\newcommand{\VrRsPreparsedAscTwentyFivePNinetyNine}{595.35}% data/runs/2026-09-14T06-41-56Z-verifyrate-rs256/verify_rate_curve.csv
\newcommand{\VrRsPreparsedAscTwoHundredCount}{36001}% data/runs/2026-09-14T06-41-56Z-verifyrate-rs256/verify_rate_curve.csv
\newcommand{\VrRsPreparsedAscTwoHundredLoad}{2.06}% data/runs/2026-09-14T06-41-56Z-verifyrate-rs256/verify_rate_curve.csv
\newcommand{\VrRsPreparsedAscTwoHundredMean}{73.92}% data/runs/2026-09-14T06-41-56Z-verifyrate-rs256/verify_rate_curve.csv
\newcommand{\VrRsPreparsedAscTwoHundredPFifty}{64.65}% data/runs/2026-09-14T06-41-56Z-verifyrate-rs256/verify_rate_curve.csv
\newcommand{\VrRsPreparsedAscTwoHundredPNinetyNine}{351.86}% data/runs/2026-09-14T06-41-56Z-verifyrate-rs256/verify_rate_curve.csv
\newcommand{\VrRsPreparsedDescFiftyCount}{9000}% data/runs/2026-09-14T06-41-56Z-verifyrate-rs256/verify_rate_curve.csv
\newcommand{\VrRsPreparsedDescFiftyLoad}{3.29}% data/runs/2026-09-14T06-41-56Z-verifyrate-rs256/verify_rate_curve.csv
\newcommand{\VrRsPreparsedDescFiftyMean}{193.82}% data/runs/2026-09-14T06-41-56Z-verifyrate-rs256/verify_rate_curve.csv
\newcommand{\VrRsPreparsedDescFiftyPFifty}{141.50}% data/runs/2026-09-14T06-41-56Z-verifyrate-rs256/verify_rate_curve.csv
\newcommand{\VrRsPreparsedDescFiftyPNinetyNine}{590.01}% data/runs/2026-09-14T06-41-56Z-verifyrate-rs256/verify_rate_curve.csv
\newcommand{\VrRsPreparsedDescFiveCount}{900}% data/runs/2026-09-14T06-41-56Z-verifyrate-rs256/verify_rate_curve.csv
\newcommand{\VrRsPreparsedDescFiveLoad}{1.87}% data/runs/2026-09-14T06-41-56Z-verifyrate-rs256/verify_rate_curve.csv
\newcommand{\VrRsPreparsedDescFiveMean}{292.11}% data/runs/2026-09-14T06-41-56Z-verifyrate-rs256/verify_rate_curve.csv
\newcommand{\VrRsPreparsedDescFivePFifty}{300.68}% data/runs/2026-09-14T06-41-56Z-verifyrate-rs256/verify_rate_curve.csv
\newcommand{\VrRsPreparsedDescFivePNinetyNine}{599.32}% data/runs/2026-09-14T06-41-56Z-verifyrate-rs256/verify_rate_curve.csv
\newcommand{\VrRsPreparsedDescFourHundredCount}{72001}% data/runs/2026-09-14T06-41-56Z-verifyrate-rs256/verify_rate_curve.csv
\newcommand{\VrRsPreparsedDescFourHundredLoad}{3.32}% data/runs/2026-09-14T06-41-56Z-verifyrate-rs256/verify_rate_curve.csv
\newcommand{\VrRsPreparsedDescFourHundredMean}{58.88}% data/runs/2026-09-14T06-41-56Z-verifyrate-rs256/verify_rate_curve.csv
\newcommand{\VrRsPreparsedDescFourHundredPFifty}{61.41}% data/runs/2026-09-14T06-41-56Z-verifyrate-rs256/verify_rate_curve.csv
\newcommand{\VrRsPreparsedDescFourHundredPNinetyNine}{149.67}% data/runs/2026-09-14T06-41-56Z-verifyrate-rs256/verify_rate_curve.csv
\newcommand{\VrRsPreparsedDescHundredCount}{18001}% data/runs/2026-09-14T06-41-56Z-verifyrate-rs256/verify_rate_curve.csv
\newcommand{\VrRsPreparsedDescHundredLoad}{4.07}% data/runs/2026-09-14T06-41-56Z-verifyrate-rs256/verify_rate_curve.csv
\newcommand{\VrRsPreparsedDescHundredMean}{111.88}% data/runs/2026-09-14T06-41-56Z-verifyrate-rs256/verify_rate_curve.csv
\newcommand{\VrRsPreparsedDescHundredPFifty}{93.64}% data/runs/2026-09-14T06-41-56Z-verifyrate-rs256/verify_rate_curve.csv
\newcommand{\VrRsPreparsedDescHundredPNinetyNine}{546.44}% data/runs/2026-09-14T06-41-56Z-verifyrate-rs256/verify_rate_curve.csv
\newcommand{\VrRsPreparsedDescTenCount}{1800}% data/runs/2026-09-14T06-41-56Z-verifyrate-rs256/verify_rate_curve.csv
\newcommand{\VrRsPreparsedDescTenLoad}{2.62}% data/runs/2026-09-14T06-41-56Z-verifyrate-rs256/verify_rate_curve.csv
\newcommand{\VrRsPreparsedDescTenMean}{286.13}% data/runs/2026-09-14T06-41-56Z-verifyrate-rs256/verify_rate_curve.csv
\newcommand{\VrRsPreparsedDescTenPFifty}{281.79}% data/runs/2026-09-14T06-41-56Z-verifyrate-rs256/verify_rate_curve.csv
\newcommand{\VrRsPreparsedDescTenPNinetyNine}{927.27}% data/runs/2026-09-14T06-41-56Z-verifyrate-rs256/verify_rate_curve.csv
\newcommand{\VrRsPreparsedDescTwentyFiveCount}{4501}% data/runs/2026-09-14T06-41-56Z-verifyrate-rs256/verify_rate_curve.csv
\newcommand{\VrRsPreparsedDescTwentyFiveLoad}{2.12}% data/runs/2026-09-14T06-41-56Z-verifyrate-rs256/verify_rate_curve.csv
\newcommand{\VrRsPreparsedDescTwentyFiveMean}{251.70}% data/runs/2026-09-14T06-41-56Z-verifyrate-rs256/verify_rate_curve.csv
\newcommand{\VrRsPreparsedDescTwentyFivePFifty}{245.76}% data/runs/2026-09-14T06-41-56Z-verifyrate-rs256/verify_rate_curve.csv
\newcommand{\VrRsPreparsedDescTwentyFivePNinetyNine}{594.08}% data/runs/2026-09-14T06-41-56Z-verifyrate-rs256/verify_rate_curve.csv
\newcommand{\VrRsPreparsedDescTwoHundredCount}{36001}% data/runs/2026-09-14T06-41-56Z-verifyrate-rs256/verify_rate_curve.csv
\newcommand{\VrRsPreparsedDescTwoHundredLoad}{3.00}% data/runs/2026-09-14T06-41-56Z-verifyrate-rs256/verify_rate_curve.csv
\newcommand{\VrRsPreparsedDescTwoHundredMean}{73.94}% data/runs/2026-09-14T06-41-56Z-verifyrate-rs256/verify_rate_curve.csv
\newcommand{\VrRsPreparsedDescTwoHundredPFifty}{64.60}% data/runs/2026-09-14T06-41-56Z-verifyrate-rs256/verify_rate_curve.csv
\newcommand{\VrRsPreparsedDescTwoHundredPNinetyNine}{349.40}% data/runs/2026-09-14T06-41-56Z-verifyrate-rs256/verify_rate_curve.csv
\newcommand{\VrSavedFifty}{631.38}% data/runs/2026-09-14T05-57-22Z-verifyrate-hs256/verify_rate_curve.csv
\newcommand{\VrSavedFive}{706.82}% data/runs/2026-09-14T05-57-22Z-verifyrate-hs256/verify_rate_curve.csv
\newcommand{\VrSavedFourHundred}{413.44}% data/runs/2026-09-14T05-57-22Z-verifyrate-hs256/verify_rate_curve.csv
\newcommand{\VrSavedHundred}{491.89}% data/runs/2026-09-14T05-57-22Z-verifyrate-hs256/verify_rate_curve.csv
\newcommand{\VrSavedTen}{741.27}% data/runs/2026-09-14T05-57-22Z-verifyrate-hs256/verify_rate_curve.csv
\newcommand{\VrSavedTwentyFive}{750.10}% data/runs/2026-09-14T05-57-22Z-verifyrate-hs256/verify_rate_curve.csv
\newcommand{\VrSavedTwoHundred}{423.23}% data/runs/2026-09-14T05-57-22Z-verifyrate-hs256/verify_rate_curve.csv


% Two columns of narrow measure with long identifiers (PSAO_JWT_KEYFORM,
% jwt.verify) produce unavoidable underfull boxes under strict justification.
\usepackage{microtype}
\emergencystretch=3em
% Long unhyphenatable identifiers are the main source of loose lines in a
% narrow two-column measure; giving TeX legal break points is the real fix.
\hyphenation{
  micro-ser-vices micro-ser-vice micro-bench-mark micro-bench-marks
  cryp-to-graph-ic cryp-tog-ra-phy au-then-ti-ca-tion va-li-da-tion
  in-stru-men-ta-tion con-fig-u-ra-tion im-ple-men-ta-tion
  Ku-ber-netes side-car side-cars name-space through-put
  pre-parsed re-pro-duc-i-ble de-ploy-ment mea-sure-ment
}
% \paragraph is used for un-numbered points, which makes hyperref see a jump in
% bookmark levels; cap the bookmark depth rather than renumber the structure.
\hypersetup{bookmarksdepth=3}

\raggedbottom

\begin{document}

\title[It Is Not the Cryptography]{It Is Not the Cryptography: Where the Cost of
Token Validation Actually Goes in Cloud-Native Microservices}

\author[1]{\fnm{Katha} \sur{Sikdar}}\email{ferdouskatha35@gmail.com}
\affil[1]{\orgdiv{Independent Researcher}}

\abstract{The runtime cost of API security in microservices is routinely
attributed to cryptographic work. We show that this attribution is wrong by two
orders of magnitude. Measuring JSON Web Token validation inside a running
Node.js service under layered Zero Trust enforcement --- edge TLS, Istio mutual
TLS and application-level validation --- we find that the signature primitive
accounts for under one percent of the per-request cost. Approximately
\VrConvPctFive\% of the cost at the arrival rates typical of published
evaluations is key material re-converted from a string on every call, an
avoidable defect that a one-line change removes. In isolation on the host, a
single HS256 validation costs \MbHsHalfKbFull~$\mu$s; inside the deployed
container the same call costs \KfHsString~$\mu$s with a string secret and
\KfHsPreparsed~$\mu$s with a pre-parsed key object, a factor of \KfHsRatio. We
further show that per-call cost depends strongly on arrival rate --- rising from
\VrHsPreparsedAscFourHundredMean~$\mu$s at \FourHundred~requests per second to
\VrHsPreparsedAscFiveMean~$\mu$s at \Five~--- and that this is neither
just-in-time compilation warmth nor garbage collection but the inter-arrival gap
itself. Because in-situ duration distributions can only be inflated by host
contention, these findings are conservative. We also report that relocating
validation into the service mesh, the remedy our own earlier work proposed,
duplicates rather than moves the work unless the application trusts a
sidecar-populated header, and that doing so opened an authentication bypass in
our testbed which we demonstrate and close. The practical consequences are that
keys must be pre-parsed, that algorithm selection must not be driven by
primitive benchmarks, and that offloading is an architectural choice with a
trust-boundary cost rather than a remedy for this overhead.}

\keywords{microservices, service mesh, JSON Web Token, performance measurement,
Zero Trust, benchmarking methodology}

\maketitle

\section{Introduction}\label{sec:intro}

A recurring claim in the literature on cloud-native security is that
authentication imposes a measurable per-request cost, and that this cost is
cryptographic. The claim is intuitive: validating a JSON Web Token involves
verifying a signature, signatures involve cryptography, and cryptography is
expensive relative to parsing. It also has an appealing consequence, namely that
the cost can be moved --- to a sidecar proxy, to an API gateway, to dedicated
hardware --- because cryptographic work is the kind of work that specialised
components do well.

We report that the claim is wrong, and wrong by a wide margin, for the most
widely deployed configuration we could measure. In a Node.js service under edge
TLS, Istio mutual TLS and application-level HS256 validation, the signature
primitive accounts for \MbHsHalfKbPrimitivePct\% of the cost of the whole
\texttt{jwt.verify()} call measured in isolation, and under one percent of the
per-request cost measured in deployment. The dominant term is not cryptography
and not parsing. It is the conversion of key material from a string into a key
object, performed once per call because the application passes a string, which
is what the library's interface invites and what published testbeds --- ours
included --- have done.

This distinction matters for three reasons. First, it changes the remedy. A cost
that is cryptographic is a candidate for offloading; a cost that is repeated
type conversion is a candidate for being deleted. Second, it changes how
algorithms should be compared. Measured on primitives alone, RS256 is far more
expensive than HS256; measured end to end in deployment, the difference is much
smaller, because a fixed per-call overhead dominates both. Third, it is a
methodological warning. The gap we report is invisible to a microbenchmark run
on a developer machine, and our own isolated measurements understated the
deployed cost by more than an order of magnitude.

\paragraph{Contributions} We make the following contributions, all supported by
measurements in the replication package~\cite{psaoartifact2026}.

\begin{enumerate}
\item An in-situ measurement of token validation inside a running service,
      using a histogram that times the verification call alone, rather than
      inferring its cost by differencing scenario averages
      (Section~\ref{sec:method}).
\item A decomposition showing that the signature primitive is a single-digit
      percentage of the isolated call and under one percent of the deployed
      cost, and that key conversion is the dominant term
      (Section~\ref{sec:results}).
\item Evidence that per-call cost depends on arrival rate, with an ascending
      and descending sweep that rules out just-in-time warmth, and a controlled
      experiment isolating the inter-arrival gap as the mechanism
      (Section~\ref{sec:results}).
\item A like-for-like algorithm comparison with both algorithms on pre-parsed
      key material, showing that primitive benchmarks mis-rank deployment cost
      (Section~\ref{sec:results}).
\item A security result: mesh-based offloading duplicates validation unless the
      application trusts a sidecar-populated header, and trusting that header
      opened an authentication bypass, which we demonstrate, map to
      API2:2023~\cite{owasp2023api}, and close (Section~\ref{sec:offload}).
\item A correction to our own earlier study~\cite{sikdar2026quantifying}, whose
      per-request attribution this paper supersedes
      (Section~\ref{sec:related}).
\end{enumerate}

\section{Background and related work}\label{sec:related}

\subsection{Layered enforcement and sidecar overhead}

Cloud-native deployments commonly stack several independent enforcement
mechanisms: TLS at the edge, mutual TLS between workloads, and application-level
token validation. Each is administered separately, and their costs are usually
reported separately.

Sidecar overhead is the best-characterised of the three. Zhu et
al.~\cite{zhu2023dissecting} decompose service mesh sidecar cost and show that
the dominant contributors depend on configuration, with inter-process
communication and socket writes dominating for TCP proxying and protocol parsing
dominating for HTTP proxying. Their central methodological point --- that the
overhead is tied to configuration and workload rather than being a fixed tax ---
is one we reach independently for validation cost. Bremler Barr et
al.~\cite{bremlerbarr2024mtls} compare mutual TLS across Istio, Istio Ambient,
Linkerd and Cilium, attributing differences to sidecar versus sidecarless
architecture. Singh et al.~\cite{singh2025resilient} evaluate Istio's
effectiveness in Kubernetes, and Hanada and Ishibashi~\cite{hanada2025slo}
address latency and failure trade-offs in microservice call chains. None of
these isolate application-level token validation, which is the layer we measure.

\subsection{Token validation cost}

Compared with mesh overhead, the cost of token validation itself is thinly
evidenced. It is frequently reported as a difference between scenario averages
--- a configuration with validation enabled minus one without --- which
attributes to validation every difference between the two runs, including
scheduling, queueing and measurement noise. Our own earlier
study~\cite{sikdar2026quantifying} did exactly this, and
Section~\ref{sec:results} shows what that approach captured.

\subsection{Zero Trust in cloud-native systems}

Zero Trust replaces perimeter trust with continuous verification of every
request. Recent systematic reviews by Gambo and Almulhem~\cite{gambo2025zta} and
Mushtaq et al.~\cite{mushtaq2025zta} survey its architecture and adoption
barriers. The relevant consequence for performance work is that verification
becomes per-request and unavoidable, so its cost is paid on every call rather
than once per session --- which is precisely why the per-call figure, and its
correct attribution, matters.

\subsection{Event-driven runtimes}

Node.js executes application logic on a single thread, so any synchronous
per-request work occupies the event loop and delays every other pending request.
Kabamba et al.~\cite{kabamba2026nodecompass} develop runtime tracing for
single-threaded environments, and Tymoshenko et
al.~\cite{tymoshenko2026predictive} argue that CPU-based autoscaling signals do
not capture event loop saturation for Node.js workloads. Peng et
al.~\cite{peng2025performance} address performance prediction and adaptive
resource adjustment for cloud-native microservices more generally. Our
instrumentation follows this line in measuring event loop delay directly, but
our claim concerns where per-call time goes rather than capacity.

\subsection{Microbenchmarks versus production behaviour}

This is the literature our central result speaks to most directly, and it was
absent from our earlier work. Schiavio et al.~\cite{schiavio2026misleading} show
that microbenchmarks executed in isolation can cause a managed runtime to
collect unrealistic execution profiles, producing aggressive optimisations that
would not occur in a real application, and therefore misleading results even
when established guidelines are followed. Traini et
al.~\cite{traini2022steady} show that steady state is frequently not reached at
all, and that developers routinely misjudge the end of the warmup phase. Our
results are an instance of exactly this hazard in a different runtime: the same
call, on the same machine, costs \TightPreparsed~$\mu$s in a tight loop and
\GapTwoHundredMsPreparsed~$\mu$s when invoked at realistic intervals.

\subsection{Cryptographic library performance}

Work on cryptographic cost in protocols concentrates on handshake and key
exchange. Montenegro et al.~\cite{montenegro2026postquantum} compare
post-quantum handshake performance across QUIC and TLS, and Souvatzidaki and
Limniotis~\cite{souvatzidaki2025pqtls} analyse post-quantum key exchange in TLS
1.3. These measure primitives and protocol exchanges carefully. Our finding is
complementary and cautionary: in an application-level validation path, the
primitive is not what dominates, so primitive-level measurement is a poor
predictor of deployed cost.

\subsection{Relationship to our earlier study}

This paper corrects~\cite{sikdar2026quantifying}. That study reported a
per-request validation cost derived by differencing amortised pod CPU between a
full-stack scenario and a mutual-TLS-only scenario, and attributed the
difference to signature verification. The present work measures the validation
call directly rather than by differencing, and finds that the attribution was
incorrect: the signature primitive is a negligible fraction of the cost, and the
dominant term is per-call key conversion. We state the correction explicitly in
Section~\ref{sec:results} and do not reuse any figure from that study.

\section{Methodology}\label{sec:method}

\subsection{Testbed}

The testbed is a single-node Kubernetes cluster (v\EnvKubernetesVersion) running
Istio with automatic sidecar injection, an injected NGINX ingress controller
terminating edge TLS, and two Node.js services. Requests traverse edge TLS, the
ingress proxy, a mutual TLS leg between sidecars, and the application. The
application under test exposes a single JSON endpoint and validates a bearer
token on every request. Node.js is v\EnvNodeVersion. Full manifests, image
digests and per-run provenance are in the replication
package~\cite{psaoartifact2026}.

The service selects its behaviour from environment variables rather than from
separate builds, so that every configuration we compare runs from one image and
one digest: \texttt{PSAO\_AUTH\_MODE} enables or disables validation,
\texttt{PSAO\_JWT\_ALG} selects HS256 or RS256, and
\texttt{PSAO\_JWT\_KEYFORM} selects whether key material is passed as a string
or pre-parsed once at startup. Comparing configurations built from different
images would place build differences inside the comparison.

\subsection{Instrumentation}

Two instruments matter here. The first is a histogram that times the
verification call alone, wrapping it with a monotonic high-resolution clock and
recording the duration under a label naming the algorithm in force. This is the
decisive instrument: it measures the call rather than inferring its cost from a
difference between scenarios. The second is an event loop delay monitor at
one-millisecond resolution, together with a request duration histogram covering
the whole handler. Metrics are exposed on a port separate from the application
port, so that scraping does not compete with the load under test.

\subsection{Isolated microbenchmark}

For comparison against in-situ measurement we decompose \texttt{jwt.verify()}
into four stages, measured on the host outside any container: the full call;
structural decoding only, which splits the token and parses the header and
payload without checking any signature; the signature primitive with the key
already parsed outside the measurement loop; and the primitive over raw buffers.
The residual --- the full call minus decoding minus primitive --- is library
dispatch, key handling, allocation and claim validation. Each stage runs
\MbHsHalfKbIterations{} iterations after a warmup, at three payload sizes, for
both HS256 and RS256.

Repeatability is reported rather than assumed. Across three passes, every
\texttt{full\_verify} figure reproduced within \MbFullVerifySpreadPct\%.
Sub-microsecond stages are less stable and we do not quote them to three
significant figures.

\subsection{In-situ rate sweep}

The central experiment drives the deployed service at a series of fixed arrival
rates --- \Five, \Ten, \TwentyFive, \Fifty, \Hundred, \TwoHundred{} and
\FourHundred{} requests per second --- for three minutes per point, and reads
the verification histogram before and after each point. Because the histograms
are cumulative, differencing the two readings isolates the distribution for that
rate alone without restarting the process.

Two design choices deserve statement. First, the process is warmed before the
first point, so that the lowest rate is not also the first traffic the process
has ever served. Second, the sweep runs ascending and then descending on the
same process. A single ascending pass confounds arrival rate with accumulated
just-in-time warmth, because the highest rate would benefit from every earlier
point's compilation; a descending pass that reproduces the ascending curve rules
that out.

\subsection{Why in-situ duration is robust to contention}

The load generator shares a host with the cluster, and host load varied across
our runs. This would be fatal to a throughput or saturation measurement, because
contention reduces achievable throughput. It is not fatal here. We measure a
duration distribution of a CPU-bound call, and contention can only \emph{inflate}
such a measurement. Every in-situ figure we report is therefore an upper bound,
and the central finding --- that deployed cost greatly exceeds isolated cost ---
is conservative under contention rather than produced by it. Host load is
recorded per point in the replication package so this can be checked.

\subsection{Key form and algorithm comparisons}

To separate the effect of key handling from the effect of the algorithm, we
measure all four combinations of algorithm and key form inside the deployed
container in a single process, with \KfHsStringIterations{} iterations each. We
then repeat the full rate sweep with HS256 on a pre-parsed key, so that the
HS256 and RS256 sweeps differ only in algorithm.

\section{Results}\label{sec:results}

\subsection{Isolated cost and its decomposition}

Measured on the host in a tight loop, a single HS256 validation of a 0.5~kB
token costs \MbHsHalfKbFull~$\mu$s. Of this, structural decoding is
\MbHsHalfKbDecode~$\mu$s and the signature primitive is
\MbHsHalfKbPrimitive~$\mu$s, leaving a residual of \MbHsHalfKbResidual~$\mu$s.
The primitive is \MbHsHalfKbPrimitivePct\% of the call and the residual is
\MbHsHalfKbResidualPct\%. For RS256 the full call is \MbRsHalfKbFull~$\mu$s with
a primitive of \MbRsHalfKbPrimitive~$\mu$s.

Even in isolation, then, the cryptographic primitive is a single-digit
percentage of HS256 validation. Figure~\ref{fig:decomposition} shows the
decomposition.

\begin{figure}[tb]
\centering
\includegraphics[width=\columnwidth]{../figures/output/fig_cost_decomposition.pdf}
\caption{Where time inside \texttt{jwt.verify()} goes, for a 0.5~kB token. The
signature primitive is a small fraction in every configuration; in the deployed
container with a string key, per-call key conversion dominates.}
\label{fig:decomposition}
\end{figure}

\subsection{In-situ cost greatly exceeds isolated cost}

Measured inside the running service at \TwoHundred{} requests per second, the
same HS256 validation costs \InsituVerifyMean~$\mu$s --- a factor of
approximately \InsituVersusMicrobenchRatio{} above the isolated figure. No
observation fell in the lowest histogram buckets; the isolated cost is not
attained at any percentile.

This is real processor time, not scheduling delay. With validation disabled and
all other work identical, application CPU at \AbRate{} requests per second was
\AbCpuNone{} millicores; with validation enabled it was \AbCpuJwt{} millicores.
The difference of \AbCpuDelta{} millicores corresponds to
\AbCpuPerRequestUs~$\mu$s of CPU per request, consistent with the wall-clock
histogram.

\subsection{Key conversion is the dominant term}

Measured inside the deployed container, with identical cryptography and an
identical token, changing only how the key reaches the library:

\begin{table}[t]
\caption{Mean cost of one validation inside the deployed container, by algorithm
and key form (\KfHsStringIterations{} iterations each).}
\label{tab:keyform}
\begin{tabular}{@{}llr@{}}
\toprule
Algorithm & Key form & Mean ($\mu$s) \\
\midrule
HS256 & string secret      & \KfHsString \\
HS256 & pre-parsed object  & \KfHsPreparsed \\
RS256 & PEM string         & \KfRsString \\
RS256 & pre-parsed object  & \KfRsPreparsed \\
\bottomrule
\end{tabular}
\end{table}

Pre-parsing reduces HS256 validation by a factor of \KfHsRatio, removing
\KfHsConvPct\% of the cost, and RS256 by a factor of \KfRsRatio. The effect is
larger for HS256 than for RS256: converting a string secret into a keyed-hash
key is more expensive here than parsing a PEM.

Across the rate sweep the same effect holds end to end. At \Five{} requests per
second, mean cost falls from \VrHsStringAscFiveMean~$\mu$s with a string secret
to \VrHsPreparsedAscFiveMean~$\mu$s with a pre-parsed key, a saving of
\VrSavedFive~$\mu$s, or \VrConvPctFive\% of the total. At \FourHundred{}
requests per second the saving is \VrConvPctFourHundred\%.
Figure~\ref{fig:keyform} shows the two curves.

\begin{figure}[tb]
\centering
\includegraphics[width=\columnwidth]{../figures/output/fig_keyform.pdf}
\caption{Validation cost with a string secret and with a pre-parsed key object,
across arrival rates. The shaded region is cost removed by a one-line change.}
\label{fig:keyform}
\end{figure}

\subsection{Cost depends on arrival rate}

Per-call cost falls monotonically with arrival rate. With a pre-parsed key it
falls from \VrHsPreparsedAscFiveMean~$\mu$s at \Five{} requests per second to
\VrHsPreparsedAscFourHundredMean~$\mu$s at \FourHundred. The descending pass
reproduces the ascending pass within \VrPassAgreementPct\% at every point,
which rules out accumulated just-in-time warmth: by the descending pass the
process had already served every higher rate, and the low-rate points remained
slow.

Garbage collection is also ruled out. Collection time rises with arrival rate
while per-call cost falls, and collection accounts for a small fraction of
verification time throughout.

The mechanism is the inter-arrival gap itself. In a single warmed process with
no HTTP, no server and no contention, inserting an idle interval between
successive calls reproduces the effect: with a pre-parsed key, cost rises from
\TightPreparsed~$\mu$s in a tight loop to \GapTwoHundredMsPreparsed~$\mu$s at a
200~ms gap, a factor of \GapEffectPreparsedRatio. Caches and branch predictors
do not survive the gap. Figure~\ref{fig:rate} plots the in-situ curves against
the isolated microbenchmark values.

\begin{figure}[tb]
\centering
\includegraphics[width=\columnwidth]{../figures/output/fig_verify_rate.pdf}
\caption{In-situ validation cost against arrival rate, with the isolated
microbenchmark values as horizontal references. Ascending and descending passes
are drawn separately; their agreement rules out just-in-time warmth.}
\label{fig:rate}
\end{figure}

\subsection{Algorithm choice looks different in deployment}

On primitives, RS256 is far more expensive than HS256: \MbRsHalfKbPrimitive~$\mu$s
against \MbHsHalfKbPrimitive~$\mu$s. In the deployed container with both
algorithms on pre-parsed keys, the whole call differs by a factor of
\KfAlgRatioPreparsed. Across the rate sweep the ratio is smaller still, around
\VrAlgRatioFive{} at \Five{} requests per second and \VrAlgRatioFourHundred{} at
\FourHundred. Figure~\ref{fig:algorithm} contrasts the two views.

\begin{figure}[tb]
\centering
\includegraphics[width=\columnwidth]{../figures/output/fig_algorithm.pdf}
\caption{RS256 against HS256, in deployment and on primitives. A primitive-level
comparison overstates the deployed difference, because fixed per-call overhead
dominates both.}
\label{fig:algorithm}
\end{figure}

\subsection{How much of observed latency is the application}

At \TwoHundred{} requests per second, mean client-observed latency was
\InsituClientMean~$\mu$s, while the server-side handler histogram recorded
\InsituHandlerMean~$\mu$s. The application accounts for approximately
\InsituAppSharePct\% of what the client sees; the remainder is TLS, the ingress,
two proxy hops, and the load generator competing for the same host. Any
attribution of per-request cost derived from client-side latency inherits that.

\subsection{What this means for the earlier attribution}

Our earlier study~\cite{sikdar2026quantifying} reported a per-request validation
cost obtained by differencing scenario CPU, and attributed it to signature
verification. The measurements here show that under one percent of validation
cost is the signature primitive, and that the large majority at the arrival
rates that study used is per-call key conversion. Its service passed a string
secret to the validation library on every request, which is the configuration
measured in the first row of Table~\ref{tab:keyform}. The magnitude it reported
is broadly consistent with what we measure for that configuration; the
attribution is not.

\section{Offloading, and the bypass}\label{sec:offload}

Because the earlier work proposed relocating validation into the sidecar, we
implemented that policy and measured it. The security findings are more
interesting than the performance ones.

\subsection{A mesh policy duplicates rather than relocates}

Applying an Istio \texttt{RequestAuthentication} causes the sidecar to validate
the token before proxying the request. It does not cause the application to stop
validating. Unless the handler is changed, validation now happens twice and no
work leaves the event loop. Relocation therefore requires the application to
trust something the sidecar tells it --- in our implementation, a header into
which the sidecar writes the verified claims --- and to skip its own signature
check when that header is present.

\subsection{The bypass}

That trust is the vulnerability. The sidecar overwrites the header on every
request, so a client cannot forge it \emph{while the policy is applied}. But the
policy is applied only transiently, when the controller decides to offload. In
every other state nothing sanitises the header. With no policy applied and no
\texttt{Authorization} header at all, a request carrying a hand-written claims
header was served with the protected payload: a complete authentication bypass,
corresponding to API2:2023 Broken Authentication~\cite{owasp2023api}.

\subsection{Closing it}

We close the bypass with an always-applied Envoy filter that strips any
client-supplied copy of the header on entry to the pod, before any validation
filter runs. It is matched on the connection manager rather than on the
validation filter, because in the state that needs protection no validation
filter exists. After the fix, a forged header alone is refused in both states,
a valid token is still served, and a valid token accompanied by a deliberately
malformed forged header is still served --- which demonstrates that the client's
value never reached the handler, since the handler fails closed on a malformed
header. A regression test in the replication package asserts all of these.

Two properties are load-bearing for this design, and both must be stated
wherever the offload is described. The mesh must overwrite the header, and
mutual TLS must prevent any client from reaching the application port without
traversing the filter chain. If mutual TLS is relaxed, the bypass reopens even
with the filter applied.

\section{Implications for practice}\label{sec:implications}

\paragraph{Pre-parse key material.} The single most effective change is to
convert keys once at startup rather than per request. In our deployment this
removed \KfHsConvPct\% of HS256 validation cost. It requires no architectural
change, no new component and no change to the trust boundary.

\paragraph{Do not select algorithms from primitive benchmarks.} A primitive
comparison ranks HS256 and RS256 by a wide margin; in deployment the gap is far
smaller, because fixed per-call overhead dominates. Algorithm choice should
weigh key distribution and trust-boundary properties, where the two differ
substantially, rather than primitive speed.

\paragraph{Treat offloading as an architectural choice, not a remedy.} Once keys
are pre-parsed, the remaining per-call cost is small. Offloading still removes
it, but it moves the trust boundary into the mesh and, as
Section~\ref{sec:offload} shows, introduces a header-trust relationship that must
be defended explicitly. That trade is worth making for reasons of policy
centralisation; it is hard to justify on the basis of this overhead alone.

\paragraph{Measure in situ.} Our isolated benchmark understated deployed cost by
more than an order of magnitude, and did so while being internally reproducible.
Instrumenting the call inside the running service cost little and was the only
measurement that exposed the error.

\section{Threats to validity}\label{sec:threats}

\paragraph{Single runtime.} All measurements are Node.js. The key-conversion
defect is a property of how this library is used, and other ecosystems may not
share it. The methodological point --- that isolated measurement mis-estimates
deployed cost --- is not runtime-specific, but our quantification of it is.

\paragraph{Single host, and contention.} Every measurement comes from one
laptop-class machine on which the load generator, ingress, sidecars and service
share cores. We record host load per run, and one sweep was discarded after
background indexing was found to have perturbed it. As argued in
Section~\ref{sec:method}, contention inflates in-situ duration rather than
deflating it, so the central comparison is conservative. Absolute figures should
not be compared against runs on dedicated hardware.

\paragraph{RS256 per-request cost is not verified in situ.} We measured RS256
inside the container and across a rate sweep, but the service was HS256-only for
the configuration our earlier study evaluated, and we do not verify any
per-request RS256 figure from that study. It remains unverified and we do not
restate it.

\paragraph{Capacity is out of scope.} We make no claim about throughput
ceilings, saturation points or autoscaling thresholds. Measurements bearing on
capacity were attempted and discarded as unreliable on this host; only
duration-distribution results, which are robust to contention in the direction
described, are reported.

\paragraph{One application shape.} The endpoint returns a small JSON payload.
Applications doing substantial per-request work would see validation as a
smaller fraction, though the absolute per-call figures should carry over.

\section{Conclusion and future work}\label{sec:conclusion}

The per-request cost of token validation in a cloud-native microservice is real
and worth attending to, but it is not cryptographic. In our deployment the
signature primitive is under one percent of it, and the dominant term is key
material converted from a string on every call --- a defect that a one-line
change removes, and one that published evaluations, including our own, have
measured without recognising. A secondary effect, the dependence of per-call
cost on arrival rate, is large enough that a benchmark run in a tight loop
understates realistic cost by more than an order of magnitude, and we isolate the
inter-arrival gap as its mechanism.

Future work should establish why per-call key conversion is so much more
expensive under virtualisation than on the host, which we did not determine;
extend the measurement to other runtimes and libraries to establish whether the
defect generalises; and verify per-request RS256 cost in situ, which remains
open.

\section*{Declarations}

\bmhead{Data and code availability} All measurements, the instrumentation, the
analysis pipeline and the scripts that produce every figure and numeric value in
this paper are available in the replication package~\cite{psaoartifact2026}.
Every number in this paper is generated from a file under \texttt{data/runs/};
none is entered by hand.

\bmhead{Competing interests} The author declares no competing interests.

\bibliography{not-the-cryptography}

\end{document}
